% EMOpticsQuality20261023
% QualityRevision20261023
% ModernLayout
\begingroup
\hypersetup{bookmarksdepth=3}
\tcbset{breakable/.default=false,after={\par\nobreak\smallskip}}
\newcommand{\ModernAnswerLabel}{{\bfseries\color{DeepBlue}解答}}
\RenewDocumentEnvironment{xesolution}{+b}{%
  \par\nobreak\smallskip\nobreak\begingroup\small
  \clubpenalties 2 10000 10000\displaywidowpenalties 2 10000 10000\noindent
  #1\par\endgroup\smallskip
}{}
\chapter{狭义相对论}

\section{相对性原理与光速不变}

狭义相对论首先用惯性参考系建立时空关系，也可在惯性坐标中描述加速粒子；它不把引力造成的时空弯曲包含在内。它建立在两条实验支持的原则上：物理定律在所有惯性系中形式相同；真空光速
$c$ 与光源和惯性观察者的运动无关。它并不是在牛顿力学（Newtonian mechanics）上额外增加一个“高速修正力”，而是改变不同惯性系之间
时空坐标的变换关系。

十九世纪末的干涉实验试图比较相互垂直方向的往返光程。图~\ref{ch31:fig:mm-lab}展示实验操作的历史情景；
分光镜把入射光分入两条近乎正交的臂，光在端镜反射后返回并合束。两臂的往返时间差会改变干涉相位；转动装置可比较不同方向的结果。
精确光路见第29章“迈克耳孙干涉仪”一节，不能从情景图中量取几何数据。
\historyfigure[0.82\linewidth]{assets/generated/michelson-morley-lab.png}{
迈克耳孙—莫雷干涉实验的复原情景示意（非史料照片）。厚重平台用于隔振，转动支撑用于比较不同方位的光程。\label{ch31:fig:mm-lab}}

设 $S'$ 相对 $S$ 沿 $x$ 轴以速度 $v$ 运动，且两原点在 $t=t'=0$ 重合。洛伦兹变换（Lorentz transformation）为
\begin{equation}
 x'=\gamma(x-vt),\qquad
 t'=\gamma\left(t-\frac{vx}{c^2}\right),
 \qquad y'=y,\quad z'=z,
\end{equation}
这里的 $t$ 不是某一处观察者“看到光到达眼睛”的时刻，而是 $S$ 系同步钟网在事件发生地点记录的坐标时间。
各系可用光信号同步：从钟A发光、在钟B立即反射、再回到A，令B反射时的读数等于A发射与接收读数的平均值。
两系采用同样的尺和钟，坐标轴平行，且 $|v|<c$。其中
\begin{equation}
 \gamma=\frac1{\sqrt{1-v^2/c^2}}.
\end{equation}
为什么时间变换含有位置 $x$？时空均匀性要求惯性坐标之间为线性变换。
$S'$ 原点满足 $x=vt$，所以可先写 $x'=a(x-vt)$、$t'=bt+dx$。
从共同原点分别向左右发光，光速不变要求
\[
 a(c-v)=c(b+dc),\qquad a(-c-v)=-c(b-dc).
\]
两式相加、相减，得到 $b=a$、$d=-av/c^2$，于是
$t'=a(t-vx/c^2)$。由相对性原理和空间反向对称性，逆变换的同一系数仍为 $a$，但 $v$ 换成 $-v$：
\[
 x=a(x'+vt')=a^2(1-v^2/c^2)x.
\]
取与 $v=0$ 连续相接的正根，便有 $a=\gamma$。这也说明不能在保留光速不变的同时继续要求 $t'=t$。

将变换式代入并展开，交叉项相消，有
\[
 c^2t'^2-x'^2=\gamma^2(1-v^2/c^2)(c^2t^2-x^2)=c^2t^2-x^2.
\]
因此它保持任意两个事件的时空间隔
\begin{equation}
 c^2\Delta t^2-\Delta x^2-\Delta y^2-\Delta z^2
\end{equation}
不变。间隔为正称类时、为零称类光、为负称类空。类时事件可由低于光速的物体连接；类光事件可由光信号连接；类空事件不能相互传递不超光速的因果信号。
低速极限 $v/c\to0$ 时，$\gamma\to1$，恢复伽利略变换（Galilean transformation）。


\subsection{先选事件，再选参考系}
洛伦兹变换变换的是同一事件的坐标。解题时先为事件命名，再列出哪一系中哪个坐标差为零。常用的三个约束不能相互替代：匀速钟的两次读数在该钟的静止惯性系满足 $\Delta x'=0$；运动杆在实验室测长要求 $\Delta t=0$；杆的固有长度要求在杆系同时测两端。
若 $\Delta x=600\,\mathrm m$、$\Delta t=1.00\,\mu\mathrm s$，取 $c\simeq3.00\times10^8\,\mathrm{m/s}$，则 $c\Delta t=300\,\mathrm m<\Delta x$，两事件类空。
使两事件同时的换系速度由 $\Delta t'=0$ 得到
\[
 v=\frac{c^2\Delta t}{\Delta x}=0.50c.
\]
再改用 $v=0.80c$，有
\[
 \Delta t'=\frac53(1.00-1.60)\,\mu\mathrm s=-1.00\,\mu\mathrm s.
\]
事件先后可以反转，但两事件没有亚光速因果联系，因此不构成因果矛盾。反过来，若两事件类时，使它们同地的速度为 $v=\Delta x/\Delta t$，其绝对值小于 $c$；此时才存在连接两事件的钟的静止系。

\subsection{光钟与时间延缓}

一只光钟由相距 $L$ 的两面平行镜组成，光脉冲在镜间往返一次记一个周期。钟的静止系中半周期为
$\Delta\tau/2=L/c$。这里固有周期 $\Delta\tau$ 对应下镜处发射A、返回C这两个事件，而不是空间分离的发射A、反射B之间的间隔。
若钟沿平行镜面、垂直于静止系光路的方向以速度 $v$ 运动，实验室看到光走斜边。横向镜间距不变，往返两段对称，所以
\begin{equation}
 \left(c\frac{\Delta t}{2}\right)^2
 =L^2+\left(v\frac{\Delta t}{2}\right)^2,
\end{equation}
从而 $\Delta t=\gamma\Delta\tau$。这不是光在实验室“走得更快”，而是光速仍为 $c$、路径变长。

\begin{center}\begin{minipage}{\linewidth}\centering
% 核心结论：同一发射反射事件在运动系形成光程直角三角形。
% 原型：纯几何对比；保留现行蓝灰风格，宽约12cm，字号10pt。
% 草图：静止 A↑B | 运动 A↗B；下方尺寸与标题分两行，L标在右侧。
\begin{tikzpicture}[font=\small,>=Latex]
 \begin{scope}
  \coordinate (A) at (1,0); \coordinate (B) at (1,2.4);
  \draw[very thick] (.2,0)--(1.8,0) (.2,2.4)--(1.8,2.4);
  \draw[vector] (A)--(B);
  \node[left=8pt] at (1,1.2) {$L=c\Delta\tau/2$};
  \node[below right=5pt] at (A) {A};\node[above right=5pt] at (B) {B};
  \node at (1,-1.45) {钟的静止系};
 \end{scope}
 \begin{scope}[xshift=5.7cm]
  \coordinate (A) at (0,0);\coordinate (B) at (1.8,2.4);\coordinate (D) at (1.8,0);
  \draw[very thick] (-.65,0)--(.65,0) (1.15,2.4)--(2.45,2.4);
  \draw[vector] (A)--(B);
  \node[above left=5pt] at (.9,1.2) {$c\Delta t/2$};
  \draw[guide] (B)--(D) (A)--(D);
  \draw[TextGray,<->] (2.85,0)--(2.85,2.4) node[midway,right=4pt] {$L$};
  \draw[TextGray,<->] (0,-.55)--(1.8,-.55) node[midway,below=3pt] {$v\Delta t/2$};
  \node[below left=4pt] at (A) {A};\node[above right=4pt] at (B) {B};
  \node at (1,-1.45) {实验室系};
 \end{scope}
\end{tikzpicture}

\figcap{横向光钟的半个周期：A发射、B反射。右图两段粗镜面分别表示对应事件时刻的位置，不是同一时刻的完整钟。几何示意未按下题数值绘制。\label{ch31:fig:light-clock}}
\end{minipage}\end{center}
图~\ref{ch31:fig:light-clock} 右侧斜边、竖边和横边组成直角三角形，直接给出时间延缓因子。

\begin{origexamplebox}{光钟的实验室周期}
光钟两镜间距 $L=0.50\,\mathrm m$，钟以 $0.80c$ 相对实验室横向运动。求钟的固有周期与实验室所见周期，并求半周期内光在实验室走过的路程。

\begin{center}\begin{minipage}{.96\linewidth}\centering
\begin{tikzpicture}[>=Latex,font=\small,every node/.style={fill=none},scale=.88]

\draw[very thick] (-.8,0)--(.8,0) (-.8,1.8)--(.8,1.8);
\draw[DeepBlue,->] (0,0)--(0,1.7);\draw[DeepBlue,->] (.2,1.8)--(.2,.1);
\draw[<->] (-1.1,0)--(-1.1,1.8) node[midway,left]{$L$};
\draw[->] (1.3,.9)--(3,.9) node[midway,above]{$v=0.80c$};\node[below] at (0,-.35){光钟自身的两面镜};

\end{tikzpicture}
\figcap{题图：两镜间距沿垂直运动方向量取，整只光钟相对实验室横向运动。\label{ch31:quality:q1}}
\end{minipage}\end{center}
\end{origexamplebox}
\begin{xesolution}
\begin{center}\begin{minipage}{.96\linewidth}\centering
\begin{tikzpicture}[>=Latex,font=\small,scale=.88]
\draw[DeepBlue,thick,->](0,0)--(2,1.5);
\draw[DeepBlue,thick,->](2,1.5)--(4,0);
\draw[dashed](0,0)--(4,0)(2,0)--(2,1.5);
\fill(0,0)circle(1.5pt)node[below]{$A$};
\fill(2,1.5)circle(1.5pt)node[above]{$B$};
\fill(4,0)circle(1.5pt)node[below]{$C$};
\draw[<->](0,-.6)--(2,-.6)node[midway,below]{$v\Delta t/2$};
\draw[dashed](2,1.5)--(4.7,1.5)(4,0)--(4.7,0);
\draw[<->](4.7,0)--(4.7,1.5)node[midway,right]{$L$};
\node[above left]at(1,.75){$c\Delta t/2$};
\end{tikzpicture}
\figcap{解图：实验室中光的往返路径。$A,B,C$ 属于不同时间的三个事件。\label{ch31:quality:a1}}
\end{minipage}\end{center}

\ModernAnswerLabel\quad 固有周期
\[
 \Delta\tau=\frac{2L}{c}=\frac{1.0\,\mathrm m}{c}\approx3.34\,\mathrm{ns}.
\]
$v=0.80c$ 时 $\gamma=1/\sqrt{1-0.8^2}=5/3$，故
\[
 \boxed{\Delta t=\gamma\Delta\tau\approx5.56\,\mathrm{ns}}.
\]
半周期内光程
\[
 \boxed{s=c\Delta t/2=\gamma L\approx0.833\,\mathrm m}.
\]
用直角三角形检查：横移 $v\Delta t/2\approx0.667\,\mathrm m$，与竖边 $0.50\,\mathrm m$ 合成的斜边确为 $0.833\,\mathrm m$。

图~\ref{ch31:quality:a1} 中，实验室半周期 $\Delta t/2$ 内，光程为三角形斜边 $c\Delta t/2$，镜面横移为 $v\Delta t/2$；两镜的横向间距 $L$ 不受纵向尺缩影响。
\end{xesolution}

\begin{center}\begin{minipage}{\linewidth}\centering
% 核心结论：时间轴为运动原点世界线，空间轴为同时线。
% 纯几何，约11x6cm，字号10pt；ct↑，x→，双向光线；P投影到直角轴。
% beta=0.4；事件P=(4,3)远离斜轴；所有标签透明。
\begin{tikzpicture}[font=\small,scale=.93,>=Latex]
 \draw[axis] (-3.5,0)--(5.6,0) node[right] {$x$};
 \draw[axis] (0,-.4)--(0,4.9) node[above] {$ct$};
 \draw[DeepBlue,dashdotted] (0,0)--(3.8,3.8);
 \draw[DeepBlue,dashdotted] (0,0)--(-3.2,3.2);
 \node[above left] at (-3.2,3.2) {$x=-ct$};
 \node[above right] at (3.8,3.8) {$x=ct$};
 \draw[Blue,thick] (0,0)--(1.8,4.5) node[above=4pt] {$ct'$};
 \draw[Blue,thick] (0,0)--(5.1,2.04) node[right=4pt] {$x'$};
 \coordinate (P) at (4,3);
 \draw[guide] (0,3)--(P)--(4,0);
 \fill (P) circle (1.6pt);\node[right=5pt] at (P) {P};
 \node[left=4pt] at (0,3) {$ct_P$};\node[below=4pt] at (4,0) {$x_P$};
 \node[below left] at (0,0) {O};
\end{tikzpicture}

\figcap{两个惯性系的时空轴。虚线将事件P投影到地面系的x、ct轴；两条点划线为向左右传播的光信号。\label{ch31:fig:spacetime-axes}}
\end{minipage}\end{center}
图~\ref{ch31:fig:spacetime-axes} 中，$ct'$ 轴是 $x'=0$ 的世界线，即 $x=vt$；$x'$ 轴是 $t'=0$ 的同时线，即 $ct=(v/c)x$。
采用相同的 $x,ct$ 长度比例尺时，世界线的横向变化率为 $\dd x/\dd(ct)=v/c$；图上通常所说的纵横斜率则是 $\dd(ct)/\dd x=c/v$（$v\ne0$）。两条光线的斜率均为 $\pm1$。
图上的欧氏距离却不是时空间隔，也不能用向斜轴作欧氏垂线的方法读取 $x',ct'$；斜坐标须按洛伦兹变换标定。

\section{固有时、时间延缓与长度收缩}

沿同一只钟世界线的两个读数事件，钟本身记录的时间称固有时 $\Delta\tau$。若该钟在 $S$ 中匀速运动，两事件满足
$\Delta x=v\Delta t$，在钟的静止系中则 $\Delta x'=0$。由间隔不变，
\[
 c^2\Delta\tau^2=c^2\Delta t^2-v^2\Delta t^2,
 \qquad \Delta\tau=\Delta t\sqrt{1-v^2/c^2},
\]
故
\begin{equation}
 \Delta t=\gamma\Delta\tau.
\end{equation}
运动钟在该惯性系坐标时间中的读数间隔更小。这里比较的是同一对事件，不能把不同位置的两只未同步时钟混入。

物体在自身静止系 $S'$ 测得的长度为固有长度 $L_0$。在地面系 $S$ 同时记录两端位置，必须取 $\Delta t=0$，令其位置差为 $L$。
对这两个测量事件，静止杆两端的坐标差仍为 $L_0$，因此
\[
 L_0=\Delta x'=\gamma(\Delta x-v\Delta t)=\gamma L,
\]
得到
\begin{equation}
 L=\frac{L_0}{\gamma}.
\end{equation}
长度收缩只发生在运动方向。它与同时性的相对性不可分：不同惯性系“同时测两端”对应不同的两件事件。

\begin{origexamplebox}{高速不稳定粒子的寿命与飞行距离}
某不稳定粒子在自身静止系平均寿命为 $\tau_0=2.2\,\mu\mathrm s$，在实验室以 $0.98c$ 运动。求实验室测得平均寿命和平均飞行距离。

\begin{center}\begin{minipage}{.96\linewidth}\centering
\begin{tikzpicture}[>=Latex,font=\small,every node/.style={fill=none},scale=.88]

\draw[->] (-.5,0)--(5,0) node[right]{$x$};\fill (0,0) circle (1.8pt) node[below]{产生事件};
\draw[DeepBlue,->] (0,.2)--(3.5,.2) node[midway,above]{$v=0.98c$};\draw[dashed] (4,-.5)--(4,.8) node[above]{衰变事件（平均位置待求）};
\node[below] at (2,-.7){粒子自身静止系平均寿命 $\tau_0$};

\end{tikzpicture}
\figcap{题图：同一粒子的产生和衰变事件，固有寿命在粒子静止系测量。\label{ch31:quality:q2}}
\end{minipage}\end{center}
\end{origexamplebox}
\begin{xesolution}
\ModernAnswerLabel\quad \[
 \gamma=\frac1{\sqrt{1-0.98^2}}\approx5.03.
\]
实验室寿命
\[
 \boxed{\tau=\gamma\tau_0\approx11.1\,\mu\mathrm s}.
\]
平均飞行距离
\[
 \boxed{\ell=v\gamma\tau_0
 =\frac{0.98c\times2.2\,\mu\mathrm s}{\sqrt{1-0.98^2}}
 \approx3.25\,\mathrm{km}}.
\]
在粒子静止系中也可得到相同相遇事件：实验室到探测位置的距离收缩为原来的 $1/\gamma$。
具体地，把产生事件记为A、衰变事件记为B；平均寿命是大量粒子的统计平均，
并非每个粒子都恰好在 $11.1\,\mu\mathrm s$ 后衰变。忽略减速时，每次飞行都有
$\Delta x=v\Delta t$，取平均后仍有 $\langle\Delta x\rangle=v\langle\Delta t\rangle$。
\begin{center}\begin{minipage}{\linewidth}\centering
% Contract: two simultaneous spatial snapshots, not two world-line diagrams.
% Same events A (creation) and B (decay); lengths refer to ground markers P,Q.
% Transparent labels, 12cm print width; no decorative backgrounds.
\begin{tikzpicture}[>=Latex,font=\small]
 \draw[->] (0,0)--(5.1,0) node[right]{$x$};
 \fill (0,0) circle(2pt) (4.4,0) circle(2pt);
 \node[below=5pt] at (0,0){P};\node[below=5pt] at (4.4,0){Q};
 \draw[DeepBlue,thick,->] (.3,.65)--(1.8,.65) node[midway,above=3pt]{$v$};
 \node[above=5pt] at (3.2,.65){粒子运动};
 \draw[<->] (0,-.65)--(4.4,-.65) node[midway,below=3pt]{$\ell=v\tau$};
 \node at (2.2,1.45){实验室系：P、Q静止};
 \begin{scope}[yshift=-3.1cm]
  \draw[->] (0,0)--(5.1,0) node[right]{$x'$};
  \fill (0,0) circle(2pt) (2.4,0) circle(2pt);
  \node[below=5pt] at (0,0){P};\node[below=5pt] at (2.4,0){Q};
  \draw[DeepBlue,thick,->] (2.4,.7)--(1.1,.7) node[midway,above=3pt]{$-v$};
  \node[above=4pt] at (4,.5){地面标记左移};
  \draw[<->] (0,-.65)--(2.4,-.65) node[midway,below=3pt]{$\ell'=\ell/\gamma=v\tau_0$};
  \node at (2.2,1.45){粒子系：粒子留在 $x'=0$};
 \end{scope}
\end{tikzpicture}

\figcap{产生时刻两系中的地面标记。P是产生位置，Q取在平均飞行距离处；长度仅作示意。\label{ch31:fig:particle-lifetime}}
\end{minipage}\end{center}
图~\ref{ch31:fig:particle-lifetime} 中，地面标记P、Q的固有间距为 $\ell$。
粒子系中它们整体向左运动，同时测得间距
$\ell'=\ell/\gamma\approx0.646\,\mathrm{km}$。
Q到达静止粒子所需时间为
\[
 t'=\frac{\ell'}v=\frac{\ell}{\gamma v}=\tau_0=2.2\,\mu\mathrm s.
\]
这样两系对同一相遇事件给出一致结果：实验室使用延长的寿命，粒子系使用收缩的地面距离。
\end{xesolution}


\subsection{加速运动中的固有时累积}
匀速公式用于整段旅程之前，要检查钟是否全程相对所选惯性系以同一速率运动，即 $\gamma$ 是否恒定。速度随时间变化时，每个足够短的时间段满足
\[
 \dd\tau=\sqrt{1-v(t)^2/c^2}\,\dd t,\qquad
 \tau=\int_{t_A}^{t_B}\sqrt{1-v(t)^2/c^2}\,\dd t.
\]
这给出理想钟沿既定世界线的累计读数；不是把全程平均速度放入 $\gamma$。只要各段之间钟的读数连续，便可逐段相加。
例如地面与飞船在A相遇，飞船以 $0.60c$ 匀速飞出3.00年（地面计时），在远端短时掉头，再以同速返回并在B重逢。忽略掉头所占时间，地面钟累计6.00年，飞船钟累计
\[
 \tau=2(3.00\,\mathrm{yr})\sqrt{1-0.60^2}=4.80\,\mathrm{yr}.
\]
A、B是同一对重逢事件，两只钟的累计读数可以直接比较。飞船去、回两段对应不同惯性系，不能把“运动钟慢”的互易说法在中途无条件拼接；差异应由各自世界线上的固有时积分计算。实际有限时长掉头还需把掉头段的积分加入。

\section{同时性的相对性与速度变换}

在 $S$ 中同时发生、相隔 $\Delta x$ 的两事件，在 $S'$ 中的时间差为
\begin{equation}
 \Delta t'=-\gamma\frac{v\Delta x}{c^2}.
\end{equation}
所以空间分离事件的同时性不是绝对的；但因果相关的类时或光速间隔事件次序不会因惯性系改变。

若物体在 $S$ 中速度为 $u_x,u_y$，先对同一段轨迹作微分：
\[
 \dd x'=\gamma(\dd x-v\dd t),\quad
 \dd y'=\dd y,\quad
 \dd t'=\gamma\left(\dd t-\frac{v\dd x}{c^2}\right)
 =\gamma(1-vu_x/c^2)\dd t.
\]
分别用位移微分除以同一个 $\dd t'$，得到
\begin{equation}
 u_x'=\frac{u_x-v}{1-u_xv/c^2},\qquad
 u_y'=\frac{u_y}{\gamma(1-u_xv/c^2)}.
\end{equation}
当 $u_x=c$ 时仍有 $u_x'=c$；低速极限恢复 $u_x'=u_x-v$。

\begin{origexamplebox}{车厢两端闪光的同时性}
在地面系 $S$ 中，一列沿 $+x$ 方向以 $v$ 运动的车厢前、后端在同一地面时刻发出闪光，两事件相距车厢在地面系的长度 $L$。求车厢系 $S'$ 中两次发光的时间差，并判断哪一端先发光。
\begin{center}\begin{minipage}{\linewidth}\centering
% 核心结论：题图给地面同时发光、地面长度L和车速，不给同时线或光相遇点。
% 纯几何，蓝灰，10pt，目标宽10cm；草图 A|------L------|B，v向右，t=0。
\begin{tikzpicture}[font=\small,>=Latex]
 \draw[thick] (0,0) rectangle (6,1.1);
 \fill[DeepBlue] (0,.55) circle(2pt);\fill[DeepBlue] (6,.55) circle(2pt);
 \node[left=5pt,align=center] at (0,.55) {后端A\\发光};
 \node[right=5pt,align=center] at (6,.55) {前端B\\发光};
 \draw[vector] (2,1.65)--(4,1.65) node[midway,above] {$v$};
 \draw[TextGray,<->] (0,-.55)--(6,-.55) node[midway,below] {$L$（地面同时测量）};
 \node at (3,.55) {地面系 $t=0$};
\end{tikzpicture}

\figcap{地面同一时刻的车厢位置及两端发光事件；箭头只给车速方向。\label{ch31:fig:train-question}}
\end{minipage}\end{center}
\end{origexamplebox}
\begin{xesolution}
\begin{center}\begin{minipage}{\linewidth}\centering\raggedright\ModernAnswerLabel\par\smallskip\centering
% 核心结论：地面同时的端点发光在车厢系不同时；光相遇点并非车中点。
% 纯几何时空图，约10x6cm，10pt；L=4,beta=1/3,ctmax=4.2。
% 后端x=ct/3，前端x=4+ct/3；光线A->P<-B，P=(2,2)。
\begin{tikzpicture}[font=\small,>=Latex]
 \def\trainL{4}\def\trainBeta{.3333333333}\def\trainT{4.2}
 \coordinate (A) at (0,0);\coordinate (B) at (\trainL,0);
 \coordinate (P) at ({\trainL/2},{\trainL/2});
 \draw[axis] (-.35,0)--(6.6,0) node[right] {$x$};
 \draw[axis] (0,0)--(0,4.9) node[above] {$ct$};
 \draw[TextGray,thick] (A)--({\trainBeta*\trainT},\trainT) node[above=6pt] {后端};
 \draw[TextGray,thick] (B)--({\trainL+\trainBeta*\trainT},\trainT) node[above=6pt] {前端};
 \draw[DeepBlue,very thick] (A)--(P)--(B);
 \draw[vector] (.65,.65)--(1.3,1.3);
 \draw[vector] (3.35,.65)--(2.7,1.3);
 \draw[Blue,dashed] (A)--(5.8,{5.8*\trainBeta});
 \node[right=5pt] at (5.8,{5.8*\trainBeta}) {$t'=0$};
 \fill (P) circle (1.6pt);\node[above=5pt] at (P) {P};
 \fill (A) circle (1.4pt);\fill (B) circle (1.4pt);
 \node[below left=3pt] at (A) {A};\node[below right=3pt] at (B) {B};
 \draw[TextGray,<->] (0,-.65)--(\trainL,-.65) node[midway,below=3pt] {$L$};
\end{tikzpicture}

\figcap{车厢两端在地面系同时发光。灰线为同向平行的车端世界线；蓝线在P相遇，但P不是运动中的车厢中点。\label{ch31:fig:simultaneity}}
\end{minipage}\end{center}


取前端事件减后端事件，则在地面系 $\Delta t=0$、$\Delta x=L>0$。洛伦兹变换给出
\[
 \Delta t'=\gamma\left(\Delta t-\frac{v\Delta x}{c^2}\right)
 =\boxed{-\gamma\frac{vL}{c^2}}.
\]
负号表示前端事件在车厢系更早。图~\ref{ch31:fig:simultaneity} 取后端发光事件A为原点、前端事件B为 $(L,0)$。
过A的车厢系同时线为 $ct=\beta x$，B在这条线下方，所以 $t'_B<t'_A=0$。
地面系中两光线为 $x=ct$ 与 $x=L-ct$，故相遇点P满足 $ct=L/2,x=L/2$；该时刻车厢中点却位于 $x=L/2+vL/(2c)$。
不能把光线相遇点、车厢中点及某一端的接收事件混作同一个事件。
$v\to0$ 或 $L\to0$ 时差异消失；$vL/c^2$ 具有时间量纲。
\end{xesolution}

\begin{center}\begin{minipage}{\linewidth}\centering
\begin{tikzpicture}[scale=.95,>=Latex]
  \draw[axis] (-.2,0)--(5.6,0) node[right]{$x$};
  \draw[axis] (0,-.2)--(0,4.3) node[above]{$y$};
  \draw[vector] (0,0)--(4.3,3.2) node[above right]{$\vect u$};
  \draw[component] (0,0)--(4.3,0) node[midway,below]{$u_x$};
  \draw[component] (4.3,0)--(4.3,3.2) node[midway,right]{$u_y$};
  \draw[vector] (.5,-.65)--(2.5,-.65) node[midway,below]{$\vect v$（换系）};
\end{tikzpicture}
\figcap{二维速度变换的分量约定示意，不按下题数值绘制。纵向、横向分量的变换分母相同，但横向还含 $1/\gamma$。\label{ch31:fig:velocity-components}}
\end{minipage}\end{center}
图~\ref{ch31:fig:velocity-components} 中换系速度只沿 $x$，所以 $u_x,u_y$ 必须分别代入不同的变换式。

\begin{origexamplebox}{横向速度分量的变换}
在实验室系中，粒子速度为 $u_x=0.60c$、$u_y=0.30c$。求相对实验室沿 $+x$ 方向以 $v=0.50c$ 运动的惯性系所测速度分量和速率。

\begin{center}\begin{minipage}{.96\linewidth}\centering
\begin{tikzpicture}[>=Latex,font=\small,every node/.style={fill=none},scale=.88]

\draw[->] (0,0)--(3,0) node[right]{$x$};\draw[->] (0,0)--(0,2) node[above]{$y$};
\draw[DeepBlue,thick,->] (0,0)--(2,1) node[right]{$\boldsymbol u=(.60c,.30c)$};
\draw[->] (.6,-.7)--(2.5,-.7) node[midway,below]{$S'\text{ 沿 }+x\text{ 以 }.50c\text{ 运动}$};

\end{tikzpicture}
\figcap{题图：实验室中的二维速度与沿 $+x$ 的参考系变换。\label{ch31:quality:q4}}
\end{minipage}\end{center}
\end{origexamplebox}
\begin{xesolution}
\ModernAnswerLabel\quad $\gamma_v=1/\sqrt{1-0.5^2}=2/\sqrt3$，共同分母
$1-u_xv/c^2=1-0.30=0.70$。故
\[
 u_x'=\frac{0.60c-0.50c}{0.70}=0.1429c,
\]
\[
 u_y'=\frac{0.30c}{(2/\sqrt3)(0.70)}=0.3712c.
\]
所以
\[
 \boxed{u'=\sqrt{u_x'^2+u_y'^2}\approx0.398c}.
\]
结果小于 $c$。若用伽利略变换而错误地令 $u_y'=u_y$，会漏掉横向分量的时间与同时性变换因子。
图~\ref{ch31:fig:velocity-components} 中的 $u_y$ 来自粒子相邻两位置的横向坐标差，
虽然 $\dd y'=\dd y$，计时却变成
\[
 \dd t'=\frac2{\sqrt3}\,(0.70)\dd t\approx0.8083\dd t.
\]
同样的横向位移除以更短的坐标时间，就有 $u_y'>u_y$。
这不是横向又施加了一个力，而是对同一段运动更换测量参考系。
\begin{center}\begin{minipage}{\linewidth}\centering
% Contract: same scale 6cm per c, components (.6,.3) and (.142857,.371154).
% Axes form mathematical 2D x/y orientation; dashed legs are components.
\begin{tikzpicture}[>=Latex,font=\small]
 \foreach \shift/\dx/\dy/\name in {0/3.6/1.8/S,6/.857143/2.226922/{S'}}{
  \begin{scope}[xshift=\shift cm]
   \draw[->] (-.2,0)--(4,0);\draw[->] (0,-.2)--(0,2.65);
   \draw[DeepBlue,very thick,->] (0,0)--(\dx,\dy);
   \draw[dashed,gray] (\dx,\dy)--(\dx,0);
   \node at (1.8,-.85){$\name$};
  \end{scope}
 }
 \node[above right=4pt] at (3.6,1.8){$\vect u$};
 \node[below=4pt] at (1.8,0){$0.60c$};
 \node[right=5pt] at (3.6,.9){$0.30c$};
 \node[above right=4pt] at (6.857143,2.226922){$\vect u'$};
 \node[below=5pt] at (6.45,0){$0.1429c$};
 \node[right=5pt] at (6.857143,1.15){$0.3712c$};
\end{tikzpicture}

\figcap{本题同一粒子在两系中的速度，左右使用同一比例尺。水平、竖直方向分别是各系的 $x,y$ 方向。\label{ch31:fig:velocity-two-frames}}
\end{minipage}\end{center}
图~\ref{ch31:fig:velocity-two-frames} 显示 $u_x'$ 减小而 $u_y'$ 增大，
同时总速率从 $\sqrt{0.60^2+0.30^2}c\approx0.671c$ 变为 $0.398c$。
\end{xesolution}

\begin{origexamplebox}{相向飞船看到的相对速度}
两飞船在地面系中分别以 $0.80c$ 向右和 $0.60c$ 向左运动。求右行飞船测得左行飞船的速度。

\begin{center}\begin{minipage}{.96\linewidth}\centering
\begin{tikzpicture}[>=Latex,font=\small,every node/.style={fill=none},scale=.88]

\draw[->] (-1,0)--(5,0) node[right]{地面系 $x$};
\fill (0,0) circle (2pt) node[below]{$A$};\fill (4,0) circle (2pt) node[below]{$B$};
\draw[DeepBlue,->] (0,.4)--(1.7,.4) node[midway,above]{$+0.80c$};\draw[DeepBlue,->] (4,.4)--(2.6,.4) node[midway,above]{$-0.60c$};

\end{tikzpicture}
\figcap{题图：两船速度先都相对地面取值，右行船 $A$ 是所求新参考系。\label{ch31:quality:q5}}
\end{minipage}\end{center}
\end{origexamplebox}
\begin{xesolution}
\begin{center}\begin{minipage}{.96\linewidth}\centering
\begin{tikzpicture}[>=Latex,font=\small,every node/.style={fill=none},scale=.88]

\draw[->] (-1,0)--(5,0) node[right]{地面系 $x$};
\fill (0,0) circle (2pt) node[below]{$A$};\fill (4,0) circle (2pt) node[below]{$B$};
\draw[DeepBlue,->] (0,.4)--(1.7,.4) node[midway,above]{$+0.80c$};\draw[DeepBlue,->] (4,.4)--(2.6,.4) node[midway,above]{$-0.60c$};

\end{tikzpicture}
\figcap{解图：先固定谁是观察者，再把另一船的有号地面速度变换过去。\label{ch31:quality:a5}}
\end{minipage}\end{center}

\ModernAnswerLabel\quad 取右行为正。地面系中左行飞船 $u_x=-0.60c$，右行飞船参考系速度参数 $v=0.80c$。所以
\[
 u_x'=\frac{-0.60c-0.80c}{1-(-0.60)(0.80)}
 =-\frac{1.40}{1.48}c\approx-0.946c.
\]
故相对速率为 $\boxed{0.946c}$，小于 $c$。直接相加得到 $1.40c$ 是伽利略速度变换在高速下的错误使用。

图~\ref{ch31:quality:a5} 中，$A$ 系由沿 $+x$ 以 $0.80c$ 的洛伦兹变换得到，故把 $B$ 的地面速度 $-0.60c$ 代入同一速度合成式；所得负号表示在 $A$ 系仍向左。
\end{xesolution}

\section{光行差与相对论多普勒效应}

平面波相位
\begin{equation}
 \Phi=\omega t-\vect k\cdot\vect r
\end{equation}
必须在同一事件给出相同相位。将逆变换 $t=\gamma(t'+vx'/c^2)$、$x=\gamma(x'+vt')$ 代入，
\[
 \Phi=\gamma(\omega-vk_x)t'
 -\gamma\left(k_x-\frac{v\omega}{c^2}\right)x'-k_yy'-k_zz'.
\]
与 $\Phi=\omega't'-\vect k'\cdot\vect r'$ 比较系数，便得到
\begin{equation}
 \omega'=\gamma(\omega-vk_x),\qquad
 k_x'=\gamma\left(k_x-\frac{v\omega}{c^2}\right),\qquad
 k_y'=k_y.
\end{equation}
对真空光波有 $|\vect k|=\omega/c$。若 $\theta$ 是 $S$ 中光线与 $+x$ 轴的夹角，则
\begin{equation}
 \cos\theta'=\frac{\cos\theta-\beta}{1-\beta\cos\theta},\qquad
 \beta=\frac vc,
\end{equation}
其中 $\theta'$ 是 $S'$ 中光线与 $+x'$ 轴的夹角；此式来自 $\cos\theta'=ck'_x/\omega'$。
频率同时满足 $f'=\gamma f(1-\beta\cos\theta)$，这两式必须使用同一套角度约定。
将 $(\omega/c,\vect k)$ 合写称为波四矢量（wave four-vector），四矢量的一般定义见本章后面的选读。

沿视线方向，令 $f_0$ 为光源静止系发射频率、$\beta>0$ 为源与观察者的相对速率。在光源系取光沿 $+x$ 传播。
观察者沿 $+x$ 远离时，代入 $\cos\theta=1$ 有
$f_{\rm rec}=\gamma f_0(1-\beta)$；观察者迎光运动则把换系速度改成 $-v$。利用
$\gamma(1\mp\beta)=\sqrt{(1\mp\beta)/(1\pm\beta)}$，得到
\begin{equation}
 f_{\rm rec}=f_0\sqrt{\frac{1-\beta}{1+\beta}},
 \qquad
 f_{\rm app}=f_0\sqrt{\frac{1+\beta}{1-\beta}},
\end{equation}
下标分别表示远离和接近。例如相对速率为 $0.60c$，远离时接收频率为 $f_0/2$，接近时为 $2f_0$。
$\beta\to0$ 时二者都回到 $f_0$。它同时含有经典传播距离效应和时间延缓，不能只用“波峰追赶”得到完整结果。




\subsection{从谱线频移反求速度}
已知光源静止系的谱线波长 $\lambda_0$，观察者测得沿视线远离的光源谱线为 $\lambda_{\rm obs}$。先取比值 $r=\lambda_{\rm obs}/\lambda_0=f_0/f_{\rm rec}$，再反解：
\[
 r^2=\frac{1+\beta}{1-\beta},\qquad
 \beta=\frac{r^2-1}{r^2+1}.
\]
例如 $\lambda_{\rm obs}=1.50\lambda_0$ 时，$\beta=5/13\approx0.385$，而把红移 $z=r-1=0.50$ 直接当作 $v/c$ 会高估速度。$z\ll1$ 时才有 $\beta\simeq z$。
这里把频移全部归因于狭义相对论中的相对运动；不含引力红移和宇宙膨胀的传播效应。


\paragraph{到达间隔包含钟慢与传播路程两部分}
在观察者惯性系中，光源沿视线以 $v>0$ 远离。两次发射在光源自身钟上相隔 $\Delta\tau=1/f_0$，在观察者系中相隔 $\Delta t=\gamma\Delta\tau$。第二次发射时光源又远离了 $v\Delta t$，第二束光还要多走这段路，所以
\[
 \Delta t_{\rm rec}=\Delta t+\frac{v\Delta t}{c}
 =\gamma(1+\beta)\Delta\tau,\qquad
 r=\frac{f_0}{f_{\rm rec}}=\gamma(1+\beta).
\]
平方并用 $\gamma^2=(1-\beta^2)^{-1}$，即可得到前面的 $r^2=(1+\beta)/(1-\beta)$。仅使用时间延缓会漏掉传播路程变化，也就无法分辨靠近与远离。
对沿视线靠近的源，可在同一约定下取 $\beta<0$，此时 $r<1$，反解式自然给出负速度；$r=1$ 给出零径向速度，$r\to\infty$ 时 $\beta\to1^-$，不会反推出超光速运动。非沿视线运动还须使用含观测角的完整多普勒关系。


\subsection{“横向”必须注明在哪个参考系}
取 $S$ 为光源静止系，$S'$ 观察者沿 $+x$ 以 $v$ 运动。两系的光线角度分别记为 $\theta,\theta'$。若光源系中横向发射，即 $\theta=90^\circ$，则
\[
 f'=\gamma f_0,\qquad \cos\theta'=-\beta.
\]
若观察者系中光线横向，即 $\theta'=90^\circ$，光行差式却要求 $\cos\theta=\beta$，从而
\[
 f'=\gamma f_0(1-\beta^2)=\frac{f_0}{\gamma}.
\]
两种频移相反，是因为所选的光线方向不同。在 $\beta=0.60$ 时，前者为 $1.25f_0$、后者为 $0.80f_0$；代入同一套角度变换即可互相核对。沿视线的接近、远离公式不能直接用于这种斜向问题。

\section{相对论动量与能量}

为了使动量守恒在所有惯性系中保持相同形式，质量为 $m$ 的粒子具有
\begin{equation}
 \vect p=\gamma m\vect v,
 \qquad E=\gamma mc^2.
\end{equation}
静止能量 $E_0=mc^2$，动能
\begin{equation}
 K=(\gamma-1)mc^2.
\end{equation}
能量动量关系为
\begin{equation}
 E^2=p^2c^2+m^2c^4.
\end{equation}
无静质量粒子满足 $E=pc$ 并以 $c$ 传播。低速展开
$K\simeq mv^2/2+3mv^4/(8c^2)+\cdots$，首项回到经典动能。

三维力仍定义为 $\vect F=\dd\vect p/\dd t$。对恒定静质量粒子的能量动量关系求导，
\[
 2E\frac{\dd E}{\dd t}=2c^2\vect p\cdot\frac{\dd\vect p}{\dd t},
 \qquad \frac{c^2\vect p}{E}=\vect v,
\]
两式相结合得到
\begin{equation}
 \frac{\dd E}{\dd t}=\vect F\cdot\vect v.
\end{equation}
因此功率关系保持熟悉的形式，但一般不能再写成 $\vect F=m\vect a$：平行和垂直于速度的加速度响应不同。守恒计算应优先使用四动量，而不是人为引入“随速度变化的质量”。
具体地，由 $\gamma=(1-v^2/c^2)^{-1/2}$ 求导：
\[
 \frac{\dd\gamma}{\dd t}=\gamma^3\frac{\vect v\cdot\vect a}{c^2},
 \quad \vect F=m\frac{\dd(\gamma\vect v)}{\dd t}
 =m\gamma\vect a+m\gamma^3\frac{\vect v\cdot\vect a}{c^2}\vect v.
\]
垂直于速度的分量中第二项为零；平行分量中用 $1+\gamma^2v^2/c^2=\gamma^2$，得到
\begin{equation}
 F_{\parallel}=\gamma^3ma_{\parallel},
 \qquad
 F_{\perp}=\gamma ma_{\perp}.
\end{equation}
低速极限两式都回到 $F=ma$；高速时继续增加平行速度尤其困难，而垂直力主要改变动量方向。

\begin{origexamplebox}{相对论动能反求速度}
粒子的动能等于其静止能量，即 $K=mc^2$。求粒子速度。

\begin{center}\begin{minipage}{.96\linewidth}\centering
\begin{tikzpicture}[>=Latex,font=\small,every node/.style={fill=none},scale=.88]

\draw[thick] (0,0) rectangle (1.5,.65);\node at (.75,.325){$mc^2$};
\draw[DeepBlue,thick] (1.5,0) rectangle (3,.65);\node at (2.25,.325){$K=mc^2$};
\node[above] at (1.5,.9){粒子总能量的两部分};\node[below] at (1.5,-.35){静质量 $m$，速度 $v$ 待求};

\end{tikzpicture}
\figcap{题图：题设给出动能与静止能相等，两部分组成粒子的总能量。\label{ch31:quality:q6}}
\end{minipage}\end{center}
\end{origexamplebox}
\begin{xesolution}
\begin{center}\begin{minipage}{.96\linewidth}\centering
\begin{tikzpicture}[>=Latex,font=\small,every node/.style={fill=none},scale=.88]

\draw[thick] (0,0) rectangle (1.5,.65);\node at (.75,.325){$mc^2$};
\draw[DeepBlue,thick] (1.5,0) rectangle (3,.65);\node at (2.25,.325){$K=mc^2$};
\node[above] at (1.5,.9){粒子总能量的两部分};\node[below] at (1.5,-.35){静质量 $m$，速度 $v$ 待求};

\end{tikzpicture}
\figcap{解图：区分静止能、动能和总能量，避免把 $K$ 直接等同于 $\gamma mc^2$。\label{ch31:quality:a6}}
\end{minipage}\end{center}

\ModernAnswerLabel\quad 由 $K=(\gamma-1)mc^2$ 得 $\gamma-1=1$，所以 $\gamma=2$。于是
\[
 \frac1{\sqrt{1-v^2/c^2}}=2,
 \qquad
 \boxed{v=\frac{\sqrt3}{2}c\approx0.866c}.
\]
若误用经典动能 $mv^2/2=mc^2$ 会得到 $v=\sqrt2c$，直接暴露出经典公式已超出适用范围。

图~\ref{ch31:quality:a6} 中，左块为静止能，右块为题给动能，总能量才等于 $\gamma mc^2$；因此先得到 $\gamma=2$，再由 $1-v^2/c^2=1/\gamma^2$ 求速率。
\end{xesolution}


\subsection{电场做功与磁场偏转的不同作用}
对于电荷量为 $q$、静质量为 $m$ 的粒子，静电场中的能量增量为 $q(\varphi_i-\varphi_f)$。电子从静止经加速电压大小 $U$ 获得 $K=eU$，所以
\[
 \gamma=1+\frac{eU}{mc^2},\quad
 v=c\sqrt{1-\left(1+\frac{eU}{mc^2}\right)^{-2}},\quad
 pc=\sqrt{eU(eU+2mc^2)}.
\]
磁场力与速度垂直，故不改变总能量。若速度垂直于匀强磁场，动量方向以角速度 $v/r$ 转动，$|\dd\vect p/\dd t|=pv/r$；与 $|q|vB$ 比较得
\[
 r=\frac p{|q|B},\qquad \omega=\frac{|q|B}{\gamma m}.
\]
高速时仍可用“动量除以电荷磁场”求半径，但不能再把动量直接写成 $mv$，也不能把回旋频率当作与能量无关的常数。低能极限两式恢复经典结果。

\section{系统不变质量、碰撞与反应阈值}
\begingroup\tcbset{breakable/.default=true,after={\par\smallskip\penalty0}}\begin{starread}{四维矢量、快度与动量中心系}
定义四维位置 $x^\mu=(ct,\vect x)$、四维动量 $p^\mu=(E/c,\vect p)$。采用度规符号
$(+,-,-,-)$ 时
\[
 p^\mu p_\mu=\frac{E^2}{c^2}-p^2=m^2c^2.
\]
能量守恒和三维动量守恒可统一写成四动量守恒。四速度 $U^\mu=\dd x^\mu/\dd\tau$ 的模恒为 $c$；
对于恒定静质量，四力 $K^\mu=\dd p^\mu/\dd\tau=m\,\dd U^\mu/\dd\tau$ 与四速度正交：对 $U^\mu U_\mu=c^2$ 求导，得到 $U_\mu\dd U^\mu/\dd\tau=0$。
四维语言的价值在于把参考系变换协变地组织起来，并不意味着时间和空间在测量意义上完全相同。
四加速度、本征力与加速杆的进一步讨论见高阶选读A6“常见四矢量与相对论动力学”。

沿同一直线的速度可用快度（rapidity）$\eta$ 表示：$\beta=\tanh\eta$、$\gamma=\cosh\eta$、$\gamma\beta=\sinh\eta$。
连续两次同向洛伦兹推促对应快度相加，所以自动给出相对论速度合成律。

% EMSpacetimeBoxSplit
\end{starread}\par\penalty0
\begin{starread}{系统不变质量与四电流}
对一个封闭粒子系，令总四动量为
\[
 P^\mu=\left(\frac{E}{c},\vect P\right).
\]
若 $P^\mu$ 为类时矢量，总动量为零的参考系称动量中心系（center-of-momentum frame），其相对实验室的速度为
\[
 \vect V_{\rm cm}=\frac{c^2\vect P}{E}.
\]
系统的不变质量 $M$ 由
\[
 M^2c^4=E^2-P^2c^2
\]
定义。在动量中心系中 $E_{\rm cm}=Mc^2$。碰撞或粒子反应的阈值应由总四动量守恒和这个不变量求出；把各粒子的静质量简单相加，通常会漏掉系统内部的相对动能与结合能。

% CompletenessEM:ch31-current-force
\noindent\textbf{四电流为什么按四矢量变换？}
先看只有一种载流子的连续介质。设其自身静止系中的有符号电荷密度为
$\rho_*$，在当前参考系中的速度为 $\vect u$。电荷不变而沿运动方向的体积缩短，
故 $\rho=\gamma_u\rho_*$、$\vect j=\rho\vect u$，从而
\[
 J^\mu=(c\rho,\vect j)=\rho_*U^\mu,
 \qquad U^\mu=\gamma_u(c,\vect u).
\]
右边是标量乘四速度，所以是四矢量；多种载流子的贡献相加，结论不变。
沿 $x$ 轴以速度 $v$ 换系时，直接使用与 $(ct,x)$ 相同的变换矩阵：
\[
 \begin{pmatrix}c\rho'\\j_x'\end{pmatrix}
 =\gamma_v\begin{pmatrix}1&-v/c\\-v/c&1\end{pmatrix}
 \begin{pmatrix}c\rho\\j_x\end{pmatrix}.
\]
因此
\[
 \boxed{\rho'=\gamma_v\left(\rho-\frac{vj_x}{c^2}\right),\qquad
 j_x'=\gamma_v(j_x-v\rho)},\qquad j_y'=j_y,\quad j_z'=j_z.
\]
连续性方程可统一写成 $\partial_\mu J^\mu=0$，其意义是同一个电荷守恒律在各系中成立，
并不要求各系在各自同时面上测得相同的局部电荷密度。

\begin{qualityfigure}
\begin{tikzpicture}[>=Latex,font=\small,scale=.85]
 \draw[thick](-.4,.3)rectangle(8,1.85);
 \draw[thick](-.4,-2.1)rectangle(8,-.55);
 \node[left]at(-.5,1.05){$S$};\node[left]at(-.5,-1.3){$S'$};
 \node[anchor=west]at(-.15,1.5){离子：$+\rho_0$，静止};
 \node[anchor=west]at(-.15,.65){电子：$-\rho_0$};
 \draw[->,DeepBlue,thick](7.5,.7)--(6.3,.7)node[midway,above]{$u$};
 \node[anchor=west]at(-.15,-.85){离子：$\gamma_v\rho_0$};
 \draw[->,DeepBlue](7.5,-.9)--(6.3,-.9)node[midway,above]{$v$};
 \node[anchor=west]at(-.15,-1.65){电子：$-\gamma_v\rho_0(1+vu/c^2)$};
 \draw[->,DeepBlue](7.5,-1.75)--(6.3,-1.75)node[midway,above]{$|u_e'|$};
 \draw[->](-.2,2.45)--(2.1,2.45)node[right]{$x$，$S'$相对$S$的速度$v$};
\end{tikzpicture}
\figcap{中性导线的两种载流子分别换系；框内标的是各系测得的体密度，不是相同同时面上的粒子计数。\label{complete:em:current-boost}}
\end{qualityfigure}

\noindent\textbf{算例：中性导线换系后为何带电？}
在图~\ref{complete:em:current-boost} 的 $S$ 系中，离子静止，密度 $+\rho_0$；
电子以 $-u\vect e_x$ 运动，密度 $-\rho_0$。所以总密度为零，但
$j_x=\rho_0u>0$。分别变换而非把两种物质当成同一个静止介质，得到
\[
 \begin{aligned}
 \rho_+'&=\gamma_v\rho_0,&j_+'&=-\gamma_vv\rho_0,\\
 \rho_e'&=-\gamma_v\rho_0(1+vu/c^2),&
 j_e'&=\gamma_v\rho_0(u+v).
 \end{aligned}
\]
电子速度也满足 $u_e'=-(u+v)/(1+vu/c^2)$，与 $j_e'=\rho_e'u_e'$ 一致。
两项相加，
\[
 \rho'=-\frac{\gamma_vv\rho_0u}{c^2},\qquad
 j_x'=\gamma_v\rho_0u.
\]
例如 $u=0.10c$、$v=0.60c$ 时 $\gamma_v=1.25$，
离子密度为 $1.25\rho_0$，电子密度为 $-1.325\rho_0$，总密度为
$-0.075\rho_0$，而 $j_x'=1.25j_x$。这里 $\rho_0$ 是两种载流子在 $S$ 系中的密度大小，
不是同时声称它们在各自静止系中也有相同固有密度。

\noindent\textbf{四力密度及其检验。}
沿用度规 $(+---)$，取电磁场张量分量
$F^{0i}=-E_i/c$、$F^{ij}=-\varepsilon_{ijk}B_k$，且
$J_\nu=(c\rho,-\vect j)$。缩并得到
\[
 f^\mu=F^{\mu\nu}J_\nu
 =\left(\frac{\vect j\cdot\vect E}{c},\,
 \rho\vect E+\vect j\times\vect B\right)
 =\left(\frac wc,\vect f\right).
\]
因张量与四矢量缩并仍是四矢量，同一矩阵作用于 $(w/c,f_x)$，所以
\[
 \boxed{w'=\gamma_v(w-vf_x),\qquad
 f_x'=\gamma_v\left(f_x-\frac{vw}{c^2}\right)},\qquad
 f_y'=f_y,\quad f_z'=f_z.
\]
用前面的中性导线做直接检验：只计外加纵向电场 $\vect E=E_x\vect e_x$、外加 $\vect B=0$
造成的贡献，则 $f_x=0$、$w=j_xE_x$。
沿电场换系后 $E_x'=E_x$、$\vect B'=0$，故由电荷密度直接算出
\[
 f_x'=\rho'E_x=-\frac{\gamma_vv j_xE_x}{c^2},\qquad
 w'=j_x'E_x=\gamma_vj_xE_x,
\]
与四矢量变换逐项相符。这个检验只隔离外场贡献，导线自身场另行叠加。
此处 $\vect f$ 是单位体积的力；它与单粒子三维力 $\vect F$ 的换系规律不同，不能混用。


\end{starread}\endgroup
\begin{origexamplebox}{对称完全非弹性碰撞后的静质量}
两个静质量均为 $m$ 的粒子在动量中心系中以相同速率 $v$ 相向运动，碰撞后结合成一个静止复合体，能量不向外辐射。求复合体静质量 $M$，并解释为何一般大于 $2m$。
\begin{center}\begin{minipage}{\linewidth}\centering
% 核心结论：已知两等质量粒子同速相向结合，待求末态静质量。
% 纯几何两行，蓝灰10pt，目标宽9cm；m -> <- m / M=?，无结果或辅助四动量。
\begin{tikzpicture}[font=\small,>=Latex]
 \node[draw,circle,minimum size=9mm] (a) at (-2.6,1) {$m$};
 \node[draw,circle,minimum size=9mm] (b) at (2.6,1) {$m$};
 \draw[vector] (a.east)--(-.5,1) node[midway,above] {$v$};
 \draw[vector] (b.west)--(.5,1) node[midway,above] {$v$};
 \node[left=8pt] at (a.west) {初态};
 \node[draw,minimum width=2cm,minimum height=.8cm] (f) at (0,-.65) {$M\;?$};
 \node[right=8pt] at (f.east) {结合后静止};
 \node[left=8pt] at (-3.05,-.65) {末态};
\end{tikzpicture}

\figcap{两粒子的初态及未知末态质量；动量中心系中初态总动量为零。\label{ch31:fig:collision-question}}
\end{minipage}\end{center}
\end{origexamplebox}
\begin{xesolution}
\begin{center}\begin{minipage}{\linewidth}\centering\raggedright\ModernAnswerLabel\par\smallskip\centering
% 核心结论：相向动量抵消，但初动能进入末态不变质量。
% 纯对比，上排初态m -> <- m，下排末态M；约10x4cm，10pt。
% 反应连接箭头与动量箭头分层，文字不骑线。
\begin{tikzpicture}[font=\small,>=Latex]
 \node[draw,circle,minimum size=10mm,inner sep=3pt] (a) at (-2.7,1.2) {$m$};
 \node[draw,circle,minimum size=10mm,inner sep=3pt] (b) at (2.7,1.2) {$m$};
 \draw[vector] (a.east)--(-.8,1.2) node[midway,above=5pt] {$+p$};
 \draw[vector] (b.west)--(.8,1.2) node[midway,above=5pt] {$-p$};
 \node[left=8pt] at (a.west) {初态};
 \node[draw,minimum width=3cm,minimum height=1cm,inner sep=6pt] (f) at (0,-1.2) {$M=2\gamma m$};
 \node[left=8pt] at (-3.2,-1.2) {末态};
 \draw[TextGray,->] (0,.65)--(f.north);
 \node[right=8pt] at (0,.25) {结合且无能量逸出};
 \node[below=6pt] at (f.south) {$\vect P=0$};
\end{tikzpicture}

\figcap{动量中心系中两个等质量粒子对撞并结合。总三动量为零，但系统静质量包含初始动能。\label{ch31:fig:inelastic-collision}}
\end{minipage}\end{center}
图~\ref{ch31:fig:inelastic-collision} 中两支初动量相消，末态复合体静止，因此能量守恒可直接求其不变质量。


初态总三动量为零，总能量为
\[
 E_i=2\gamma mc^2.
\]
末态复合体静止，能量为 $Mc^2$。能量守恒给出
\[
 \boxed{M=2\gamma m}.
\]
当 $v>0$ 时 $\gamma>1$，所以 $M>2m$；增加的静质量来自被困在复合体内部的初始动能（内能）。
对应的能量等式为 $(M-2m)c^2=2(\gamma-1)mc^2$：左侧是静能增加，右侧是初始总动能。
$v\to0$ 时 $M\to2m$。若碰撞中有光子或其他粒子带走能量，必须把其四动量计入，不能仍用此式。
\end{xesolution}





\begin{origexamplebox}{固定靶粒子反应的阈能}
静质量为 $m_a>0$ 的入射粒子撞击实验室中静止、静质量为 $m_b>0$ 的靶粒子。反应末态至少有两个有质量粒子，其静质量之和为 $M_f>m_a+m_b$，即反应为吸能反应。
忽略外力，只讨论由能量和动量守恒决定的运动学阈值，不考虑库仑势垒及反应概率。求最小入射动能，并说明放能反应为什么不能直接使用所得负值作为动能下限。

\begin{center}\begin{minipage}{.96\linewidth}\centering
\begin{tikzpicture}[>=Latex,font=\small,every node/.style={fill=none},scale=.88]

\fill (0,0) circle (2pt) node[below]{$a,m_a$};\draw[DeepBlue,->] (0,.3)--(1.7,.3) node[midway,above]{$\boldsymbol p_a$};
\fill (2.5,0) circle (3pt) node[below]{$b,m_b,\ v_b=0$};
\draw[->] (3.1,0)--(4.3,0);\draw[dashed] (5,0) circle (.5);\node[above] at (5,.7){末态总静质量至少 $M_f$};

\end{tikzpicture}
\figcap{题图：实验室固定靶反应的初态；入射总动量非零，末态在实验室不能全部静止。\label{ch31:quality:q8}}
\end{minipage}\end{center}
\end{origexamplebox}
\begin{xesolution}
\ModernAnswerLabel\quad 先在动量中心系中寻找生成末态所需的最少能量。在这个系中，总三动量为零，而每个有质量粒子的能量满足
\[
 E_i^*=\sqrt{m_i^2c^4+p_i^{*2}c^2}\ge m_ic^2.
\]
相加得 $E_{\rm final}^*\ge\sum_i m_ic^2=M_fc^2$。
只有所有 $p_i^*=0$ 时取等号，这就是阈值末态“没有相对动能”的由来。
注意它们只在质心系共同静止：实验室初态总动量不为零，末态必须整体运动来承接这部分动量，不能在实验室把末态全画成静止。
\begin{center}\begin{minipage}{\linewidth}\centering
% Contract: generic two-mass final state at threshold. All lab momenta forward;
% CM final particles at rest. Directions only; lengths not magnitudes.
\begin{tikzpicture}[>=Latex,font=\small,y=.75cm]
 \node[anchor=east] at (0,.8){实验室初态};
 \fill (.7,.8) circle(2.5pt) (3.7,.8) circle(2.5pt);
 \node[below=5pt] at (.7,.8){$a$};\node[below=5pt] at (3.7,.8){$b$ 静止};
 \draw[DeepBlue,thick,->] (.9,.8)--(2.5,.8) node[midway,above=4pt]{$p_a$};
 \node[anchor=east] at (0,-.9){实验室阈值末态};
 \foreach \x in {.7,3.7}{\fill (\x,-.9) circle(2.5pt);\draw[DeepBlue,thick,->] (\x,-.9)--++(.9,0);}
 \node[above=5pt] at (.7,-.9){$1$};\node[above=5pt] at (3.7,-.9){$2$};
 \node at (2.2,-1.6){同一速度 $V_c$，整体运动};
 \node[anchor=east] at (0,-2.8){质心系阈值末态};
 \fill (.7,-2.8) circle(2.5pt) (3.7,-2.8) circle(2.5pt);
 \node[below=5pt] at (.7,-2.8){$p_1^*=0$};\node[below=5pt] at (3.7,-2.8){$p_2^*=0$};
\end{tikzpicture}

\figcap{以两个末态粒子为例的阈值几何，箭头只表示动量方向。\label{ch31:fig:threshold-frames}}
\end{minipage}\end{center}
图~\ref{ch31:fig:threshold-frames} 第二行各末态粒子具有同一实验室速度 $V_c$，
但质量不同则动量不同；第三行才是所有动量均为零的质心系。
因此阈值末态总四动量的不变质量为 $M_f$。设入射粒子的实验室总能量、动量为 $E_a,p_a$。初态总能量和总动量为
\[
 E_i=E_a+m_bc^2,
 \qquad P_i=p_a.
\]
利用 $E_a^2-p_a^2c^2=m_a^2c^4$，初态不变量为
\begin{align}
 E_i^2-P_i^2c^2
 &=m_a^2c^4+m_b^2c^4+2m_bc^2E_a.
\end{align}
在阈值处令它等于 $M_f^2c^4$，得到
\[
 E_{a,\rm th}=\frac{M_f^2-m_a^2-m_b^2}{2m_b}c^2.
\]
入射动能是总能量减静止能量，所以
\begin{equation}
 \boxed{K_{\rm th}=\frac{M_f^2-(m_a+m_b)^2}{2m_b}c^2}.
\end{equation}
若末态只是原粒子的弹性散射，$M_f=m_a+m_b$，阈能为零；若产生了新的静质量，阈能不仅要提供静能差，还要保证实验室系总动量守恒，因此通常大于简单的静能差。量纲上括号内是质量，乘 $c^2$ 后为能量。
若改成具有连续末态相空间的放能反应，$M_f<m_a+m_b$，入射动能趋于零时初态就已有足够的不变能量，释放的静能差成为末态相对动能，不能再令末态相对动能为零。
因此这类反应仅从运动学考虑的下限为0。把两种情形合写为
\[
 K_{\min}=\max\left\{0,\frac{M_f^2-(m_a+m_b)^2}{2m_b}c^2\right\}.
\]
例如 $m_a=m_b=m$、$M_f=1.5m$ 时，代数表达式为 $-0.875mc^2$，只说明反应放能，不表示可以有负入射动能。
单一粒子末态属于另外的能量匹配问题，不能把上面的连续末态阈值讨论直接套用。
\end{xesolution}

\subsection{用光子系统检验不变质量}
单个真空光子的静质量为零，并不意味着多个光子组成的系统不变质量也为零。两光子的能量分别为 $E_1,E_2$，传播方向夹角为 $\theta$，总动量模平方为
\[
 P^2c^2=E_1^2+E_2^2+2E_1E_2\cos\theta.
\]
由总能量 $E=E_1+E_2$，得到
\[
 M^2c^4=E^2-P^2c^2=2E_1E_2(1-\cos\theta).
\]
同向时 $M=0$，不存在总动量为零的惯性系；反向且等能 $E_1=E_2=\varepsilon$ 时，总动量为零，$M=2\varepsilon/c^2$。这里的质量属于系统，不是给每个光子分配了静质量。计算碰撞阈值时使用总四动量，正是为了保留这类内部相对运动的能量。


\hypersetup{bookmarksdepth=2}
\endgroup
